To specify the affine identifications and the gluing, use the
saturated basis
\[
(1,0,0,0),\qquad(0,1,1,0),\qquad(0,1,0,1)
\]
of $\ker R$, with splitting vector $w=v_1$.  In these coordinates put
\[
\begin{gathered}
a_0=(-1,0,0),\quad a_1=(3,1,1),\\
b_0=(-3,-1,0),\quad b_1=(-3,0,-1),\quad b_2=(-1,0,0),\quad b_3=0.
\end{gathered}
\]
The two summands in the negative slice are the convex hulls of
the images $A(v)=\varphi_0(v)+v_2+w$ and $B(v)=\varphi_1(v)$,
with the tail cone added.  The vertex images, in the order
$v_2,\ldots,v_8$, are
\[
A:\ (a_0,a_1,a_1,a_0,a_0,a_1,a_1),\qquad
B:\ (b_3,b_1,b_0,b_0,b_1,b_3,b_2).
\]
The tail rays are
\[
\begin{aligned}
q_0&=(-4,-1,0),&q_1&=(-4,0,-1),&q_2&=(0,0,1),\\
q_3&=(0,1,0),&q_4&=(1,0,0),&q_5&=(2,1,1),&q_6&=(3,1,1).
\end{aligned}
\]
Table~\ref{tab:x9slices} lists all maximal original charts.  An
index set in the $A$ or $B$ column means the convex hull of those
vertices plus the indicated tail.  The marking records complete
locus.  Write $A_j,B_j$ for the summands in row $j$.
All remaining cells are obtained from the nonzero face
cones of the twelve listed facets by the same slice construction.

\begin{table}[ht]
\centering
\caption{The two displaced slices and their markings.}
\label{tab:x9slices}
\begin{tabular}{rccccc}
\toprule
facet & vertices & tail rays & $A$ & $B$ & marked \\
\midrule
$0$ & $0136$ & $134$ & $01$ & $1$ & yes \\
$1$ & $0137$ & $346$ & $1$ & $13$ & yes \\
$2$ & $0145$ & $024$ & $01$ & $0$ & yes \\
$3$ & $0147$ & $246$ & $1$ & $03$ & yes \\
$4$ & $0156$ & $014$ & $0$ & $01$ & yes \\
$5$ & $02367$ & $4$ & $01$ & $13$ & no \\
$6$ & $02457$ & $4$ & $01$ & $03$ & no \\
$7$ & $0256$ & $4$ & $0$ & $013$ & no \\
$8$ & $134568$ & $01235$ & $01$ & $012$ & yes \\
$9$ & $1378$ & $356$ & $1$ & $123$ & yes \\
$10$ & $1478$ & $256$ & $1$ & $023$ & yes \\
$11$ & $2345678$ & $\varnothing$ & $01$ & $0123$ & no \\
\bottomrule
\end{tabular}
\end{table}
